\section{Воспроизводимость}\label{sec:repro}

Весь код, точные сертификаты и данные открыты:
\begin{center}
\url{https://github.com/ivaleo/Chromatic}\quad (лицензия MIT).
\end{center}
Для пяти утверждений таблицы~\ref{tab:main} достаточно одного маршрута
каждому (таблица~\ref{tab:manifest}); все пути --- относительно каталога
\path{audit-data/}, команды выполняются из него.
Единый запуск
\begin{center}
\texttt{python -m chromatic\_research.campaigns.verify\_main\_results}
\end{center}
пересчитывает итоговые неравенства всех пяти строк из опубликованных
дробей на чистой рациональной арифметике за доли секунды (уровень~1 ниже),
а с флагом \texttt{--full} запускает и полные верификаторы строк $4$, $5$,
$7$ и $9$ (последний перечисляет $1\,654\,230$ вершин и занимает несколько
часов).

\begin{table}[htbp]\centering\footnotesize
\caption{Пять утверждений: вход, верификатор, ожидаемый вывод, уровень
независимости. Уровни: 1~--- пересчёт итоговых неравенств из опубликованных
дробей; 2~--- другой алгоритм в том же пакете; 3~--- независимая реализация
без общих геометрических процедур; 4~--- аналитическое доказательство.
Файлы данных лежат в \texttt{results/}, верификаторы --- в
\texttt{chromatic\_research/campaigns/}, тесты --- в \texttt{tests/}.}
\label{tab:manifest}
\setlength{\tabcolsep}{4pt}
\renewcommand{\arraystretch}{1.15}
\begin{tabular}{@{}l >{\raggedright\arraybackslash}p{4.4cm} >{\raggedright\arraybackslash}p{4.2cm} >{\raggedright\arraybackslash}p{3.5cm} c@{}}
\toprule
утверждение & вход & верификатор & ожидаемый вывод & ур.\\
\midrule
$\chi(\R^4)\le43$ & \path{r4_k43_eisenstein_rational.json} & \path{dim4_k43_verify.py} & $d^2=\frac{298768\ldots}{296327\ldots}>\ell_0^2$ & 3\\
 & \path{cert43_eisenstein.json} & опорный сертификат & \texttt{ok}; $D^2_{\downarrow}\ge\ell_0^2\diam^2P$ & 2\\
$\chi(\R^5)\le132$ & \path{metric_deform_a5_132_refined_certificate.json} & \path{certificate_window_fix.py} & $38$ векторов окна, запас $>0$ & 2\\
 & \path{metric_deform_a5_132_refined_independent_exact_audit.json} & \path{verify_exact_voronoi.py} & те же $R^2$, $D_{\min}^2$ & 3\\
$\chi(\R^7)\le1029$ & \path{dim7_1029_exact.json} & \path{dim7_1029_exact.py} & $D_{\min}^2=7$; \texttt{interval\_valid} & 2\\
 & --- & \path{dim7_1029_layer_cert.py} & $\max\varphi<7/4$ & 3\\
$\chi(\R^9)\le7203$ & \path{dim9_7203_exact.json} & \path{dim9_7203_exact.py} & $D_{\min}^2=7$; \texttt{interval\_valid} & 2\\
 & следствие~\ref{cor:dim9} ($9604$) & \path{test_product_calculus.py} & $6/7+1/9=61/63<1$ & 4\\
$\chi(\R^{10})\le45619$ & параметры двух блоков & \path{test_product_calculus.py} & $6/7+4/31=214/217<1$ & 4\\
\bottomrule
\end{tabular}
\end{table}

Каждый сертификат самодостаточен: рациональная форма Грама, матрица
подрешётки, точный квадрат радиуса покрытия, полный список векторов окна и
KKT-сертификат проекции для каждого из них; проверить его может короткая
программа, не имеющая ничего общего с кодом, которым он получен. Уровень~3
для строк $4$, $5$ и $7$ означает независимость \emph{кода}: перепроверки
выполнены программами, написанными с нуля и не разделяющими с основным
конвейером геометрических процедур (для $1029$ второй путь --- слоёный
сертификат замечания~\ref{rem:r7second}), но не независимыми исполнителями в смысле слепой репликации. Неизменяемая версия: настоящий текст отвечает коммиту
репозитория, указанному в файле \path{paper/article1/README.md};
архив с DOI будет указан при подаче.

\paragraph{Вклад авторов и использование ИИ.}
Постановка задачи, метод и вычислительный конвейер, конструкция индекса
$132$, ламинирование, продуктовое исчисление с оценками $9604$ и $45619$ и
планарная граница --- Л.\,Л.~Иванов. Н.\,В.~Глушковой принадлежат решётка
индекса~$43$ (первая точка семейства~\eqref{eq:eisform}, подсказавшая
симметрийное сужение) и точные рациональные верификаторы ламинирований
$E_6^*/343$ и $E_8/2401$, переведшие $1029$ и $7203$ из численных результатов
в теоремы~\ref{thm:r7ex} и~\ref{thm:r9ex}; во втором случае решающей была
точная унимодулярная замена базиса, без которой перечисление вершин
девятимерной ячейки не доводится до конца. Симметрийное сужение поиска и
независимые перепроверки --- совместные. Постатейный учёт с датами и коммитами ведётся
в репозитории (\path{CONTRIBUTIONS.md}). Вычислительный конвейер, скрипты
кампаний и первоначальные редакции текста первого автора создавались с
систематическим использованием ИИ-ассистента (большие языковые модели
общего назначения в агентной оболочке; коммиты несут трейлеры
\texttt{Co-Authored-By}); постановки, направление поиска, критерии отбора и
итоговая проверка принадлежат авторам, и для проверяемости теорем
происхождение поискового кода несущественно: каждая из них сведена к
конечному списку неравенств между явно выписанными рациональными числами.
Ответственность за все утверждения несут авторы.
